\chapter{Endovascular Aneurysm Repair (EVAR)}

\section*{Introduction \& Clinical Indications}
Endovascular aneurysm repair (EVAR) treats a suitable abdominal aortic or iliac artery aneurysm by placing a covered stent-graft inside the vessel. The graft excludes the aneurysm from arterial pressure while preserving flow through the aorta and iliac arteries. EVAR may be considered for symptomatic or ruptured aneurysms, or for intact aneurysms whose rupture risk outweighs the procedural risk. Repair thresholds depend on diameter, growth, symptoms, sex, anatomy, life expectancy, and shared decision-making; there is no single threshold for every patient.

EVAR is an alternative to open repair, not a universally superior replacement. It usually has a smaller initial physiologic burden, but requires suitable landing zones and lifelong imaging surveillance because endoleak, graft migration, limb occlusion, and sac enlargement can occur. Complex aneurysms may require fenestrated, branched, or iliac devices, or open surgery.

\section*{Relevant Surgical Anatomy}
Planning evaluates the aortic neck, renal and visceral arteries, aneurysm sac, iliac bifurcations, internal iliac arteries, access vessels, and collateral circulation. The diameter, length, angulation, thrombus, calcification, and tortuosity of the proximal and distal seal zones determine device feasibility. Femoral access may be percutaneous or surgical. Preservation of at least one internal iliac artery is usually prioritized when possible to reduce pelvic ischemic complications.

\section*{Preoperative Workup \& Preparation}
Thin-slice CT angiography with appropriate arterial-phase reconstruction is used for sizing and planning. Assessment includes renal function, anemia, cardiopulmonary risk, frailty, contrast allergy, anticoagulants, infection, and the possibility of conversion to open repair. Ruptured aneurysm requires rapid resuscitation, permissive hypotension when appropriate, blood-product preparation, and multidisciplinary coordination. Consent covers endoleak, renal injury, limb ischemia, bowel or spinal ischemia, access complications, conversion, reintervention, and death.

\section*{Surgical Technique \& Procedure}
Through femoral access, guidewires and delivery sheaths are advanced under fluoroscopy. The main graft is positioned below or across the intended proximal landing zone, and iliac limbs are deployed to exclude the sac while maintaining branch-vessel perfusion. Balloon molding may improve apposition. Completion angiography assesses renal arteries, endoleak, limb patency, and pelvic flow. Access sites are closed and distal pulses are documented.

\section*{Complications \& Risks}
Risks include access-site bleeding or pseudoaneurysm, arterial dissection, limb thrombosis or kinking, endoleak, graft migration, renal injury, contrast reaction, distal embolization, infection, bowel or spinal ischemia, buttock claudication, erectile dysfunction, rupture, and need for conversion or secondary intervention. Ruptured aneurysm carries substantially higher risk than elective repair.

\section*{Postoperative Care \& Recovery}
Monitoring includes blood pressure, urine output, renal function, hemoglobin, abdominal symptoms, groin wounds, and lower-limb pulses. Early mobilization and antithrombotic therapy are individualized. Follow-up imaging is essential to detect endoleak, sac enlargement, graft integrity, and limb complications. A new abdominal or back pain, syncope, cold or painful leg, bleeding, fever, or reduced urine output requires urgent assessment.

\section*{Outcomes \& Prognosis}
EVAR can provide durable aneurysm exclusion in appropriately selected anatomy, but durability depends on seal zones, device, anatomy, surveillance, and reintervention. Long-term outcome is influenced by cardiovascular disease, smoking, renal function, and aneurysm-related complications. EVAR should be chosen through anatomy-based comparison with open repair rather than a universal claim of superiority.

\section*{References}
\begin{enumerate}
  \item Society for Vascular Surgery. Guidelines for the Care of Patients With an Abdominal Aortic Aneurysm.
  \item European Society for Vascular Surgery. Clinical Practice Guidelines on the Management of Abdominal Aorto-Iliac Artery Aneurysms.
  \item Society for Vascular Surgery. Perioperative Care Resource for Vascular Disease.
\end{enumerate}

\clearpage
